\begin{table}[htbp]
\centering
\footnotesize
\setlength{\tabcolsep}{4pt}
\caption{更新型比值样本的二项稀释与含州截距拟合。面板 A 报告每个保留概率 $q$ 下原始规模 $C_t$ 上的斜率、隐含稀释斜率 $c_q=[b(q)-b(1)]/(1/q-1)$、截距变化 $\Delta a$ 以及 $C_t\ge50$ 子样本（$n=828$）中的 $c_q$。面板 B 报告 $1/C_t$ 上的斜率以及各设定下的斜率截距比。括号内数值为点估计及 95\% 州$\times$周双向自助法百分位区间。面板 B 中的区间来自各规格共享的一套 1000 次抽取，因此合并区间与表~\ref{tab:s1} 略有微小出入。}
\label{tab:s11}
\begin{tabular}{lllll}
\toprule
\multicolumn{5}{l}{\textit{A. 二项稀释（$n=1748$）}}\\
$q$ & $b(q)$ & $c_q$ & $\Delta a$（$\times10^{-3}$） & $c_q$（$C_t\ge50$） \\
\midrule
0.8 & 0.678 [0.447, 0.918] & 0.474 [0.038, 0.956] & $-$0.34 [$-$1.84, 1.17] & 0.472 [$-$0.509, 1.416] \\
0.6 & 0.881 [0.609, 1.197] & 0.481 [0.189, 0.883] & $-$0.89 [$-$4.13, 1.46] & 0.441 [$-$0.123, 1.071] \\
0.5 & 1.053 [0.728, 1.397] & 0.493 [0.230, 0.818] & $-$1.30 [$-$4.78, 1.84] & 0.458 [$-$0.032, 0.988] \\
0.4 & 1.316 [0.917, 1.788] & 0.504 [0.264, 0.785] & $-$1.90 [$-$6.26, 1.87] & 0.451 [0.074, 0.946] \\
0.3 & 1.731 [1.256, 2.371] & 0.502 [0.319, 0.765] & $-$2.40 [$-$7.48, 1.54] & 0.458 [0.086, 0.827] \\
0.2 & 2.457 [1.736, 3.181] & 0.474 [0.307, 0.652] & $-$2.50 [$-$8.50, 3.39] & 0.457 [0.195, 0.773] \\
0.1 & 4.010 [2.913, 5.346] & 0.383 [0.265, 0.532] & 0.44 [$-$7.01, 11.93] & 0.483 [0.241, 0.857] \\
\midrule
合并斜率 $c$ & -- & 0.407 [0.311, 0.531] & -- & 0.477 [0.286, 0.769] \\
泊松基准系数 $b_P$ & -- & 0.446 & -- & 0.420 \\
\bottomrule
\end{tabular}

\bigskip
\begin{tabular}{lll}
\toprule
\multicolumn{3}{l}{\textit{B. 州固定效应模型（$n=1748$，51 个州与特区）}}\\
设定规格 & 斜率 $b$ & 表观比值 $b/a$ \\
\midrule
合并横截面基准 & 0.560 [0.368, 0.771] & 37.4 [18.6, 73.9] \\
州截距，均值 $\bar{a}_s$ & 0.379 [0.104, 0.669] & 17.3 [3.3, 51.4] \\
州截距，中位数 $a_s$ & -- & 21.7 [3.7, 60.6] \\
州截距，行数加权 $\bar{a}_s$ & -- & 18.4 [3.5, 49.9] \\
州截距，病例量加权 $\bar{a}_s$ & -- & 24.6 [4.9, 58.4] \\
州截距，以 $D_t$ 为规模 & 0.369 [0.111, 0.643] & 17.7 [3.2, 51.4] \\
仅包含周固定效应截距 & 0.508 [0.232, 0.835] & -- \\
同时包含州与周两组截距 & 0.78 [$-$1.20, 2.77] & -- \\
配对比值（横截面 / 均值 $\bar{a}_s$） & -- & 2.16 [0.93, 7.75] \\
\bottomrule
\end{tabular}
\end{table}
